% Source: sources/returns/S6_04_MAPS_MANIFOLDS/code/manifolds_corrected_initial_proof.tex:1–142
\paragraph{Theorem D.2.52 (Proved: all height-integral initial pairs with corrected original additive caps).}
For any of the twenty-one recorded Mordell--Weil rank-one
profiles, take $P=aG,Q=bG$ with $a,b>0$ and
$d_0\mid ab$. Choose any nonnegative whole-fibre
linearization weights of total $ab/d_0$, supported
at multiplicative or additional regular fibres.
Use the actual corrected line bundle of D.2.50.
The following filling choices have an exact all-order
fundamental-group presentation and complete integral
homology algorithm:
\begin{enumerate}[label=(\roman*)]
\item every marked same-fan rank-one or triangular
Mumford cap at a multiplicative or regular site;
\item every free finite-good additive logarithmic cap,
with all its affine powers and all five integral Thom
columns retained;
\item every original additive cap whose corrected
multidegree has $h=1$ in (D.IC2), and every unchanged
zero-degree original additive cap;
\item every original multiplicative cap with its
uncorrected degree vector, as in D.2.46;
\item arbitrarily many regular logarithmic caps and
any ordinary integral clutch.
\end{enumerate}
The odd-star fibres are filled originally in this
algorithm. In particular, the recorded Smith
corrections from D.2.50 admit every positive pair
with integral height: all newly non-narrow additive
positions are the component pair $(2,2)$ at an odd
star, and their corrected degrees are (D.IC8).
Newly non-narrow multiplicative positions may all
use the marked Mumford filling.
Thus this theorem removes the narrowness-OR restriction
on the initial sections in D.2.46.
It does not identify original multiplicative caps
with their nonzero vertical corrections, nor the
unresolved marked $h>1$ additive caps.

The exact implementation is
\path{code/manifolds_corrected_initial.py}.
It accepts the same unbounded data as D.2.46 and
optional literal component-divisor vectors. A supplied
vector is checked against the integral equation
(D.VC3); no rational correction is substituted.
For an original additive site it computes the
orbit lengths and $h$, and applies the complete
marked kernel (D.IC7) only when $h=1$.
A nonzero whole-fibre change at a zero-degree
original cap is also retained as outside this
implementation until its relative meridian marking
is supplied.

\begin{proof}
D.2.50 constructs the actual initial line bundle
and proves that its punctured monodromy is exactly
(D.GI2), without the narrowness-OR assumption.
All entries are integral and their ordered product
is the identity. The common punctured torus bundle
therefore has exactly the same integral fibre-and-graph
bar model as in D.2.41 and D.2.46, including its
quadratic comparison and literal last-boundary word.

Every selected Mumford cap depends on the actual
marked unipotent period matrix and the selected fan.
The rank-one calculation and D.2.43 therefore apply
to these newly corrected inputs verbatim: the full
nonsaturated angular-area kernel, primitive normal
discs and all component meridians remain as stated.
No extension of the vertical divisor over a different
original cap is being inferred from this fact.
The finite-good and regular logarithmic caps use
the actual affine quotient action, so their freeness
criterion, norm columns and fifth Thom generator
are unchanged as well.

At an original additive site with $h=1$, D.2.51
supplies every required integral marking. Its
fundamental group is generated by the primitive
long circle; all three vertical directions are
killed. The fibre degree-two kernel is the actual
primitive area kernel (D.IC7), and both literal
fixed-normal tori bound. Its component contribution
to $H_4$ is the number $\rho$ of component orbits,
not the number of components before the permutation.
The unchanged zero-degree case and the uncorrected
original multiplicative case are already proved
in D.2.31, D.2.33 and D.2.46.

The implementation therefore adds these exact
local columns to the common bar presentation.
Its fundamental-group presentation retains the
commuting fibre generators, the conjugation matrices,
the actual ordered peripheral product, each finite
power relation, and each nonfinite killed vector.
The normal closure of the vertical vectors is
computed by repeatedly adjoining their conjugates
and taking integral Hermite normal form until the
lattice stabilizes. The ascending-chain condition
on submodules of $\mathbb Z^3$ proves termination.
No saturation or assumed abelian quotient replaces
this presentation.

For degree-two homology, form the full integer
relation matrix $M$ and retain its Smith cokernel.
If $E$ denotes the common piece after the finite
caps have been filled, the remaining relative
free rank is exactly
\[
           \sum_i b_4(U_i)-b_1(E)+b_1(X).
                                                        \tag{D.CI1}
\]
The proof is the same integral pair sequence as
(D.GI6): each $H_2(U_i,\partial U_i)=H^4(U_i)$
is free, and its actual boundary columns are
already included. The cap torsion permitted by
D.2.43 remains in its nonsaturated fibre kernel;
the new original additive caps have no degree-two
torsion because $h=1$. The quotient in (D.CI1)
is free, so it splits as an abelian-group extension.
This gives $H_2(X)$ with all torsion.
The Euler number is twice the sum of the indices
of the positive rank-two caps. The exact $H_1$,
this Euler number, integral duality and universal
coefficients give every remaining group, including
the Euler-zero cases.

Finally, the canonical correction list in D.2.50
has $45$ additional residue pairs over the full
profile list. At additive sites the component
orders are prime except for the two odd stars.
The criterion $\operatorname{ord}(\Phi)\mid pq$
therefore forces $p=0$ or $q=0$ except for
$p=q=2$ in $\mathbb Z/4$.
The displayed canonical correction at those
stars has a fixed inner coefficient $-1$,
so $h=1$. This proves that the algorithm applies
to every positive height-integral initial pair
with these retained corrections and the stated
filling choices. The conclusion is unbounded in
section multiples, weights, logarithmic orders
and every permitted affine twist.
\end{proof}
